% ===================================================================
\section{Discussion}
% ===================================================================

\subsection{Beyond Compliance: \DPP\ as Business Enabler}

A recurring theme at \DPPforEU~2026 was the need to ensure that \DPPs\ are perceived and implemented as more than a regulatory compliance burden. As a conference speaker emphasised at \DPPforEU~2026~\cite{dpp4eu2026_day3}, the goal is to ``\textit{make sure that the implementation and the perception of \DPPs\ stay beyond compliance}''. The portal-backed architecture supports this transition: by hosting \DPPs\ as queryable knowledge graphs in a centralised triple-store, the portal enables not only compliance checking but also advanced use cases --- supply chain analytics, circular economy business model innovation, AI-driven material recovery optimisation, and cross-product substance tracking --- that go far beyond the minimum regulatory requirements. The conference observation that \DPPs\ should ``\textit{transit from statistic documentation to the intelligent decision support system}''~\cite{dpp4eu2026_day3} is directly enabled by the portal's \SPARQL-queryable knowledge graph architecture, which transforms static \DPP\ records into a living, queryable substrate for circular economy intelligence.

\subsection{The Governance Challenge: Human Agreement, Not Technology}

A striking post-conference reflection from the DigiPass \CSA\ project crystallised a conclusion that emerged across the event: the \DPP\ challenge is fundamentally one of governance and human agreement, not technology. As Peter Klein, the DigiPass \CSA\ project leader, wrote in the post-conference summary~\cite{dpp4eu2026_postconf}:

\begin{quote}\emph{
``All the excellent presentations, posters, and podium discussions at the \DPPforEU\ crystallized one particular conclusion in my mind: the solution of the \DPP\ challenge is more related to human agreement, commitment, and collaboration than to any technical implementation issue.''
}
\end{quote}

This observation was echoed by Paul Breiner of the TRUSTex project, who described the conference's value as ``\textit{the transdisciplinary exchange of technical, regulatory, governance and industry-related perspectives around a common objective: unlocking the Digital Product Passport's potential to drive the systemic transition towards a Circular Economy}''~\cite{dpp4eu2026_postconf}. The portal-backed architecture proposed in this paper is fundamentally a governance architecture: it provides the structural framework within which human agreement, commitment, and collaboration can be operationalised. The portal is not a technological endpoint; it is the institutional infrastructure that makes cross-organisational data governance possible at scale.

\subsection{The Role of DigiPass Europe as Community Sustainer}

The conference revealed a clear aspiration within the DigiPass consortium to sustain the community beyond the \CSA's lifetime. As a speaker noted at \DPPforEU~2026~\cite{dpp4eu2026_day2}, the plan is to ``\textit{go from DigiPass to DigiPass Europe}'' in order to ``\textit{create a community around these topics that sustain the effort that has been made so far with everyone}''. The proposed portal-backed architecture aligns with this aspiration: DigiPass Europe could serve as the organisational entity that operates or oversees the federated portal infrastructure, the public repository for ontologised regulations, and the \SHACL\ brick library. The conference observation that ``\textit{we need actually for this to act a kind of central hub connecting the industry the stakeholders}''~\cite{dpp4eu2026_day2} maps directly onto the portal's role as the semantic and governance hub of the \DPP\ ecosystem. The \COST\ action model, described as a bottom-up initiative open to all interested participants~\cite{dpp4eu2026_day2}, provides a complementary funding and governance mechanism for the ongoing community work, particularly for the domain-specific semantic information generation that the \SHACL\ brick library requires.

The conference format itself --- a single-track, ``one-for-all'' programme deliberately designed to avoid community fragmentation --- reflects the same architectural principle underlying the portal-backed model. As Heinz Adolf Preisig observed in the post-conference summary: ``\textit{One could argue for separate sessions and longer talks, but it would split the community. Given the status of the \DPP\ development affair, I felt a need to see the complete picture and minimise the formation of specialised groups}''~\cite{dpp4eu2026_postconf}. The architectural parallel is direct: just as the conference community resisted fragmentation through a unified
